\textbf{1--7. Boundary conditions and eigenvalue equation.} A free end carries no axial force, so its stress is zero. With constant $E$, the two conditions are $U'(0)=U'(L)=0$. Using the course-note harmonic solution,
\[
u(x,t)=\left(C_1\sin\gamma x+C_2\cos\gamma x\right)e^{-i\omega t},
\qquad \sigma=E\gamma\left(C_1\cos\gamma x-C_2\sin\gamma x\right)e^{-i\omega t},
\]
the left condition gives $C_1=0$ for $\gamma>0$. The right condition gives $C_2\sin(\gamma L)=0$. The trivial choice $C_2=0$ gives zero vibration. A nonzero amplitude requires
\[
\boxed{\sin z=0,\qquad z=\gamma L.}
\]
The two end conditions can also be written before eliminating the amplitudes as
\[
\gamma C_1=0,\qquad
\gamma\left[C_1\cos(\gamma L)-C_2\sin(\gamma L)\right]=0.
\]
For each positive eigenvalue, the first equation removes the sine component and the second selects the cosine shape. The zero-frequency case is treated separately because division by $\gamma$ is then invalid.
\textbf{8--17. Natural frequencies and wavelengths.} For positive-frequency modes, $z_j=j\pi$ and $\gamma_j=j\pi/L$, where $j=1,2,\ldots$. Since $c=\omega/\gamma$ and $\lambda=2\pi/\gamma$,
\[
\boxed{\omega_j=\frac{j\pi c}{L},\qquad f_j=\frac{jc}{2L},\qquad \lambda_j=\frac{2L}{j},\qquad L=j\frac{\lambda_j}{2}.}
\]
The fundamental is the lowest positive natural frequency $f_1=c/(2L)$. Higher modes are overtones. Thus an integer number of half-wavelengths fits the bar.

\textbf{18--21. Displacement and stress modes.} Setting the arbitrary $C_2=1$ gives $U_j(x)=\cos(j\pi x/L)$. The course-note orthonormality convention for these positive-frequency modes is
\[
\boxed{\frac{2}{L}\int_0^L U_j(x)U_k(x)\,dx=\delta_{jk},
\qquad \Sigma_j(x)=E U'_j(x)=-\frac{Ej\pi}{L}\sin\frac{j\pi x}{L}.}
\]
The zero-frequency rigid translation $U_0=1$ is also admissible. Its stress is zero and its normalization is $L^{-1}\int_0^L U_0^2\,dx=1$. The derivation and mode convention follow the course notes, axial vibration section 3.3, PDF pages 231--235.

For $j=k$, the integral of $\cos^2(j\pi x/L)$ is $L/2$. For different positive integers, the cosine products integrate to zero. Multiplication by $2/L$ therefore produces the Kronecker delta shown above. Differentiating the displacement gives the stress shape, which vanishes at both free ends.
